%% ==========================================================================
%% DRAFT -- August 5, 2026. Author byline: Daniel Kirtchakov (05oz,
%% halfounce.io; no institutional affiliation).
%%
%% Result referee-confirmed 2026-08-05; novelty sweep clean 2026-08-05
%% (Open Problem Garden fetched live, arXiv API, web sweeps, Yuster's
%% recent output; no determination of integral nu_3(9) or nu_3(10) found
%% anywhere since Kabiya--Yuster 2008).
%%
%% Framing constraints imposed by the adversarial referee and honored
%% throughout this note:
%%   * The NEW content is the two integral LOWER bounds, established over
%%     every tournament isomorphism class (191,536 at n=9; 9,733,056 at
%%     n=10). The matching UPPER bounds are Yuster 2004's (Section 4:
%%     oriented Turan blow-up; fractional strengthening in Kabiya--Yuster
%%     2008, Section 2.4), and are credited so at every point of use. We
%%     write "no certified determination existed", never "the values were
%%     unknown".
%%   * Declared trust assumption: gentourng (nauty 2.9.3) enumerates
%%     completely; cross-validated exhaustively for n <= 7 and
%%     count-validated against OEIS A000568 and exact Burnside numbers
%%     at n = 9, 10.
%% Kabiya--Yuster read in full 2026-08-05 from an archived copy
%% (primary-source-kabiya-yuster-tt3.pdf; author preprint; retained in
%% the author's private working tree, not redistributed). Yuster 2004
%% checked verbatim against arXiv math/0304180 the same day (archived
%% copy likewise private). Fix pass 2026-08-05: see FIXLOG.md.
%% ==========================================================================
\documentclass[11pt]{amsart}

\usepackage[margin=1.15in]{geometry}
\usepackage{amsmath,amssymb}
\usepackage{booktabs}
\usepackage{array}
\usepackage[colorlinks=true,allcolors=blue]{hyperref}

\newcolumntype{L}[1]{>{\raggedright\arraybackslash}p{#1}}
\emergencystretch=3em
\tolerance=1200
\hbadness=4000

\newcommand{\TT}{\mathrm{TT}_3}
\newcommand{\nut}{\nu_3}
\newcommand{\nuf}{\nu_3^{*}}

\theoremstyle{plain}
\newtheorem{theorem}{Theorem}[section]
\newtheorem{lemma}[theorem]{Lemma}
\newtheorem{proposition}[theorem]{Proposition}
\newtheorem{corollary}[theorem]{Corollary}
\theoremstyle{definition}
\newtheorem{definition}[theorem]{Definition}
\theoremstyle{remark}
\newtheorem{remark}[theorem]{Remark}

\begin{document}

\title[Certified values of $\nut(9)$ and $\nut(10)$]{The minimum number of
arc-disjoint transitive triples in a tournament: a first certified
determination of $\nut(9)=9$ and $\nut(10)=12$}

\author{Daniel Kirtchakov}
\address{Independent researcher (\texttt{05oz}); no institutional
affiliation.}
\email{daniel@halfounce.io}
\urladdr{https://halfounce.io}

\thanks{\emph{Computation and authorship.} All searches, encodings, and
certificate designs in this work were produced by \textbf{Claude Fable~5}
(Anthropic), directed by the author, on a single Apple~M4 laptop. The
independent verifiers import only the Python standard library and share no
code with the search pipeline; the adversarial referee of
\S\ref{ssec:referee} was a separate agent instance given no access to the
pipeline and wrote its own code throughout. This is a factual methods
statement, and it is part of the point of the note: the artifacts are
designed so that the provenance of the \emph{search} is irrelevant to the
validity of the \emph{result}. External tools used by the pipeline only:
\texttt{gentourng} (nauty 2.9.3), CaDiCaL~3.0.1
(\texttt{--lrat --no-binary}), Apple clang 17.0.0, Python~3.}

\thanks{\emph{Prior-art record.} The Kabiya--Yuster text \cite{KY08} was
read in full on August~5, 2026, from a copy retained in the author's working
tree and not redistributed here (an
author preprint; section references follow its numbering). The
concluding-remarks account in \cite{Yuster04} of the verified range and of
the upper-bound construction was checked verbatim against the arXiv text
(\texttt{math/0304180}) the same day; the results of \cite{Yuster04} are
otherwise used as reported in \cite{KY08}. The novelty sweep of the same
date --- the Open Problem Garden entry fetched live, the arXiv API, and
general web sweeps --- found no determination of the integral values
$\nut(9)$ or $\nut(10)$ anywhere in the record since 2008. The upper-bound
half of both values was never open; see \S\ref{ssec:credit}. \emph{Draft
of August 5, 2026.}}

\subjclass[2020]{Primary 05C20; Secondary 05C70, 68V15}
\keywords{Tournament, transitive triple, packing, Yuster's conjecture,
certified computation, SAT, LRAT}

\date{August 5, 2026}

\begin{abstract}
For a tournament $T$ let $\nut(T)$ be the maximum number of pairwise
arc-disjoint transitive triples in $T$, and let
$\nut(n)=\min_T \nut(T)$ over all $n$-vertex tournaments. Yuster (2004)
conjectured $\nut(n)=\lceil n(n-1)/6-n/3\rceil$ and verified this for all
$n\le 8$ ($n=8$ by computer, the smaller cases by direct argument),
without published certificates; no determination of $\nut(9)$ or
$\nut(10)$, certified or otherwise, appears in the record since. We supply the first certified determination:
$\nut(9)=9$ and $\nut(10)=12$, the conjectured values.

The new content is the two lower bounds. An exact branch-and-bound packer
was run over every tournament isomorphism class --- all $191{,}536$ at
$n=9$ and all $9{,}733{,}056$ at $n=10$, enumerated by \texttt{gentourng}
with counts validated against OEIS A000568 and independently recomputed
Burnside numbers --- and emitted, for every class, an explicit packing of
$9$ (resp.\ $12$) arc-disjoint transitive triples, each re-verified from
the raw arc data by a standard-library checker sharing no code with the
search. The matching upper bounds are not new: they are Yuster's (2004),
via oriented Tur\'an blow-ups, restated with a fractional strengthening
by Kabiya and Yuster (2008), \S2.4. What we add on that side is an
artifact: for an explicit minimizer in each order ---
$C_3[C_3,C_3,C_3]$ at $n=9$ and $C_3[TT_4,C_3,C_3]$ at $n=10$ --- an
LRAT-certified proof that no packing of $10$ (resp.\ $13$) triples
exists, replayed by an independent checker that regenerates the CNF from
its specification.

One trust assumption is declared rather than checked: that
\texttt{gentourng} (nauty) enumerates the isomorphism classes completely.
It was cross-validated exhaustively for $n\le 7$ and count-validated at
$n=9,10$. Everything else was verified twice: by stdlib-only verifiers,
and by an adversarial referee who re-ran the enumeration to byte-level
set equality, re-solved both minimizers with an independent exact solver,
and exercised both verifiers with negative controls. The next open case,
$\nut(11)$ (conjectured value $15$), is a factor of roughly $93$ larger
and is within reach of the same pipeline (CPU-days, not CPU-years); it
has not been started.
\end{abstract}

\maketitle

\section{Introduction}\label{sec:intro}

\subsection{The problem}
A tournament is an orientation of the complete graph. A transitive triple
($\TT$) is a $3$-vertex subtournament with no directed cycle; a
$\TT$-packing of a tournament $T$ is a set of pairwise arc-disjoint
transitive triples of $T$, and $\nut(T)$ is the maximum size of one.
Following Kabiya and Yuster \cite{KY08} we write
\[
\nut(n)\;=\;\min\{\nut(T): T \text{ a tournament on } n
\text{ vertices}\},
\]
the value of the packing problem on the \emph{worst} $n$-vertex
tournament. Yuster \cite{Yuster04} proved that
$\nut(n)\ge\frac{51}{392}n^{2}(1-o(1))$, conjectured the exact value, and
verified the conjecture for all $n\le 8$ --- $n=8$ by computer, the
smaller cases by direct argument \cite[\S4]{Yuster04}:

\begin{center}
\emph{Conjecture (Yuster \cite{Yuster04}; Conjecture 1.1 of
\cite{KY08}).} $\ \nut(n)=\lceil n(n-1)/6-n/3\rceil$.
\end{center}

\noindent The conjectured sequence begins, from $n=3$:
$0,1,2,3,5,7,9,12,15,\dots$ Kabiya and Yuster \cite{KY08} improved the
asymptotic lower bound to $\frac{41}{300}n^{2}(1-o(1))$, still the best
known, by solving the \emph{fractional} relaxation at $r=10$ by linear
programming ($\nuf(10)=12$) and transferring the bound to the integral
problem with a loss. Their paper is the last progress recorded on the
problem's Open Problem Garden entry \cite{OPG}, which records no
progress beyond 2008.

The state of the integral problem at the start of this work was
therefore: conjectured value known for all $n$; verified for $n\le8$ by
Yuster ($n=8$ by computer), with no certificates for that range in
\cite{Yuster04} or anywhere else in the record; at $n=9$
and $n=10$, an upper-bound half that was closed in print (see
\S\ref{ssec:credit}) and a lower-bound half with nothing in the record.
\emph{No certified determination of $\nut(9)$ or $\nut(10)$ existed.}

\subsection{Result}

\begin{theorem}\label{thm:main}
$\nut(9)=9$ and $\nut(10)=12$. Consequently, together with the verified
range $n\le 8$ (\cite{Yuster04}, reproduced in \S\ref{ssec:packer}),
Yuster's conjectured formula $\lceil n(n-1)/6-n/3\rceil$ holds for all
$n\le 10$, and a counterexample, if one exists, has $n\ge 11$.
\end{theorem}

Every claim in Theorem~\ref{thm:main} is backed by a machine-checkable
artifact: an explicitly verified witness packing for every one of the
$191{,}536+9{,}733{,}056$ tournament isomorphism classes (the lower
bounds), and LRAT unsatisfiability certificates on explicit minimizer
tournaments (the upper bounds), all replayable with standard-library
Python and shipped with hashes (\S\ref{sec:artifacts}). This extends the verified
range of the conjecture from $n\le8$ (uncertified) to
$n\le10$ (certified).

\subsection{Provenance and credit}\label{ssec:credit}
The division of labor between this note and the literature is exact, and
stating it precisely is a referee-imposed condition on the writing.

\emph{The upper bounds are Yuster's.} Section 4 of \cite{Yuster04}
orients the complete $3$-partite Tur\'an graph
$V_1\to V_2\to V_3\to V_1$, completes the classes arbitrarily, and
observes that every transitive triple of the resulting tournament uses at
least one within-class arc; hence no $\TT$-packing exceeds the number of
within-class arcs, which is exactly $\lceil n(n-1)/6-n/3\rceil$. Kabiya
and Yuster \cite[\S2.4]{KY08} restate the construction --- crediting the
tournaments to \cite{Yuster04} --- and strengthen the bound to the
\emph{fractional} packing number. At $n=9$
with parts $(3,3,3)$ this gives $\nut(9)\le 9$; at $n=10$ with parts
$(4,3,3)$ it gives $\nut(10)\le 12$. So the upper-bound half of
Theorem~\ref{thm:main} was never open, and nothing about the
\emph{values} $9$ and $12$ being upper bounds is claimed here. What this
note adds on the upper side is a certificate: an LRAT proof, for an
explicit minimizer in each order, that one more triple does not fit
(\S\ref{ssec:upper}).

\emph{The lower bounds are the new content.} $\nut(9)\ge 9$ and
$\nut(10)\ge 12$ assert something about \emph{every} tournament of the
given order, and for these no argument, computation, or certificate
existed in the record. They are established here by exhaustive
per-isomorphism-class computation with per-class witnesses
(\S\ref{ssec:lower}).

\emph{The enumeration is nauty's; the counts are classical.} Tournament
isomorphism classes were generated by \texttt{gentourng} from the nauty
suite of McKay and Piperno \cite{nauty}, and the class counts
$191{,}536$ and $9{,}733{,}056$ are OEIS sequence A000568 \cite{OEIS},
recomputed independently during verification via Burnside's lemma
(Davis's formula \cite{Davis54}).

\subsection{What is not claimed}
We do not claim the upper bounds, which are \cite{Yuster04}'s (with the
fractional strengthening of \cite{KY08}). We do not claim any asymptotic
improvement; the best general lower bound remains
$\frac{41}{300}n^{2}(1-o(1))$ of \cite{KY08}. We do not claim the values
were unexpected: they are exactly the conjectured ones. We do not claim
the conjecture; it remains open for $n\ge11$. And one link in the chain
is trusted rather than machine-checked --- the completeness of the
\texttt{gentourng} enumeration --- and \S\ref{sec:tcb} states exactly
what supports that trust.

\section{The two halves of the computation}\label{sec:comp}

\subsection{Enumeration}\label{ssec:enum}
All tournaments were handled up to isomorphism, which suffices because
$\nut$ is an isomorphism invariant. \texttt{gentourng -q 9} produced
$191{,}536$ tournaments and \texttt{gentourng -q 10 $r$/16},
$r=0,\dots,15$, produced $9{,}733{,}056$ in sixteen slices (per-slice
counts are shipped in \texttt{n10\_chunk\_counts.txt} and sum exactly).
Both totals equal A000568 \cite{OEIS}. Each tournament is a line of
$\binom{n}{2}$ characters, the upper triangle row by row, \texttt{1} at
pair $(i,j)$, $i<j$, iff the arc is $i\to j$; the output strings were
checked to be pairwise distinct (\texttt{sort -u}: $191{,}536$ and
$9{,}733{,}056$).

\subsection{The exact packer, validated on the known
range}\label{ssec:packer}
\texttt{tt3pack.c} ($136$ lines of C) computes $\nut(T)$ exactly by
branch and bound: it branches on a live uncovered arc with the fewest
free triples through it --- either some free triple through it joins the
packing, or the arc is marked permanently uncovered --- and prunes with
the bound $|\text{packing}| + \lfloor\text{live free arcs}/3\rfloor$.
With a target $t>0$ it stops early as soon as a packing of size $\ge t$
is found and prints it; with target $0$ it returns the exact maximum with
an optimal packing. Before any new claim was attempted, the packer was
run in exact mode over the complete enumeration for $n=4,\dots,8$ and
returned minima $1,2,3,5,7$ --- reproducing Yuster's verified
range \cite{Yuster04} exactly (achieving tournaments and optimal packings
in \texttt{smalln\_table.txt}).

\subsection{Lower bounds: a witness for every isomorphism
class}\label{ssec:lower}
The sweep $\texttt{gentourng -q 9}\mid\texttt{tt3pack 9 9}$ terminated
with every one of the $191{,}536$ classes reporting a packing of $\ge9$
arc-disjoint transitive triples, each line carrying the packing itself
(about a second of wall clock; log \texttt{n9\_sweep.txt.gz}). The sixteen
$n=10$ slices, run as
$\texttt{gentourng -q 10 } r/16\mid\texttt{tt3pack 10 12}$, terminated
with all $9{,}733{,}056$ classes reporting a packing of $\ge12$, about
$3$\,s per slice (logs \texttt{n10\_sweep\_0..15.txt.gz}). Since every
tournament is isomorphic to a listed one and $\nut$ is invariant,
\[
\nut(9)\ge 9, \qquad \nut(10)\ge 12 .
\]
These $9{,}924{,}592$ witness lines are the note's entire novel
mathematical content, and they are what the verification of
\S\ref{sec:verif} re-checks line by line.

\subsection{Upper bounds: certified optimality on explicit
minimizers}\label{ssec:upper}
Following \cite[\S4]{Yuster04} and \cite[\S2.4]{KY08}, the candidate
minimizers are the cyclic blow-ups $C_3[T_1,T_2,T_3]$: three vertex classes with all arcs
$V_1\to V_2\to V_3\to V_1$ and arbitrary tournaments inside the classes.
The packer, in exact mode, was run on every within-class variant: all
$8$ variants with parts $(3,3,3)$ at $n=9$ have $\nut(T)=9$ exactly, and
all $48$ variants with parts $(4,3,3)$ at $n=10$ have $\nut(T)=12$
exactly (\texttt{cand9.out}, \texttt{cand10.out}) --- consistent with
\cite[\S4]{Yuster04} and \cite[\S2.4]{KY08}, where the upper bound is
proved for an arbitrary within-class completion. The representatives fixed for certification are
\[
T_9=C_3[C_3,C_3,C_3], \qquad T_{10}=C_3[TT_4,C_3,C_3]
\]
(arc strings in \texttt{minimizer9.bits}, \texttt{minimizer10.bits};
$T_{10}$ has the transitive $4$-tournament in its $4$-class). For each,
two artifacts are shipped:

\begin{itemize}
\item an optimal packing --- $9$ triples in $T_9$, $12$ in $T_{10}$
(\texttt{minimizer9.witness}, \texttt{minimizer10.witness}) --- so the
values are attained;
\item a proof that no larger packing exists: a CNF asserting ``$T$
contains $10$ (resp.\ $13$) pairwise arc-disjoint transitive triples'',
refuted by CaDiCaL~3.0.1 \cite{CaDiCaL} with an LRAT proof
\cite{LRAT}: \texttt{min9\_ge10.cnf} ($2{,}430$ variables, $5{,}058$
clauses; proof \texttt{min9\_ge10.lrat}, $419{,}396$ bytes) and
\texttt{min10\_ge13.cnf} ($5{,}740$ variables, $11{,}916$ clauses; proof
\texttt{min10\_ge13.lrat}, $1{,}938{,}465$ bytes).
\end{itemize}

\noindent The encoding (full byte-exact specification in the
\texttt{encode\_cnf.py} docstring) has one variable per transitive triple
of $T$, a binary conflict clause for each pair of triples sharing an arc,
and a Sinz sequential at-most-$K$ counter \cite{Sinz05} on the negated
triple variables, $K=m-k$ with $m$ the number of transitive triples, so
that at least $k$ triples must be selected. Its correctness splits
asymmetrically. That any satisfying assignment yields $k$ arc-disjoint
triples is forced clause by clause (conflict clauses; counter soundness).
The direction that makes UNSAT meaningful --- any packing of size $k$
extends to a satisfying assignment of the counter variables --- is the
standard completeness property of the sequential-counter encoding
\cite{Sinz05}, and it sits in the trusted base
(\S\ref{sec:tcb}). As positive controls, the same encoding at the
attained value is satisfiable, and was solved as such:
\texttt{min9\_ge9.cnf} and \texttt{min10\_ge12.cnf} are SAT (CaDiCaL exit
$10$), so unsatisfiability appears exactly at value${}+1$. Hence
$\nut(T_9)=9$ and $\nut(T_{10})=12$, giving $\nut(9)\le9$ and
$\nut(10)\le12$ --- the bounds already available on paper from
\cite[\S4]{Yuster04}, now in checkable form. With \S\ref{ssec:lower},
Theorem~\ref{thm:main} follows.

\begin{remark}[The fractional parameter]\label{rem:frac}
The blow-up bound holds for the fractional packing number $\nuf$ as
well --- this strengthening is Kabiya--Yuster's \cite[\S2.4]{KY08} ---
so
$\nuf(9)\le9$ and $\nuf(10)\le12$; combining with
$\nuf(n)\ge\nut(n)$ and Theorem~\ref{thm:main} gives
$\nuf(9)=9$ and $\nuf(10)=12$. The $r=10$ value agrees with the LP
computation reported in \cite[\S3]{KY08} and provides an independent
derivation of it that does not rest on an uncertified LP solve. This is
consistent with the observation of \cite[\S2.4]{KY08} that their
Conjecture 1.1 implies $\nuf(n)=\nut(n)$ for all $n$.
\end{remark}

\section{Verification}\label{sec:verif}

\subsection{Independent stdlib-only verifiers}\label{ssec:stdlib}
Two verifiers, sharing no code with the searcher and importing only the
Python standard library, re-check everything the pipeline produced.

\texttt{verify\_sweep.py} re-checks every one of the $9{,}924{,}592$
sweep lines from the raw arc string alone: each witness triple is
transitive in the stated tournament, no arc is used twice, the count
meets the target, the line totals match the expected class counts
exactly. All $17$ witness files pass with zero failures. A below-target
line would be flagged loudly as a candidate counterexample; there were
none.

\texttt{verify\_minimizer.py} re-establishes $\nut(T_9)=9$ and
$\nut(T_{10})=12$ from raw artifacts: it re-checks the optimal packing
against the bits; \emph{regenerates} the CNF from the encoding
specification with independent code and demands equality with the
shipped CNF as a clause multiset; and replays the LRAT proof with its
own RUP-with-hints checker --- no external tool, and RAT steps are
refused outright, which incidentally establishes that both proofs are
pure RUP. Both certificates verify.

\subsection{The adversarial-referee protocol}\label{ssec:referee}
The result was then subjected to an adversarial referee --- a separate
agent instance with no access to the pipeline, writing entirely fresh
code --- whose protocol is part of the contribution and is recorded in
the program's results ledger of 2026-08-05. The referee re-ran the
enumeration from scratch and confirmed byte-level set equality with the
shipped sweeps; recomputed the class counts $191{,}536$, $9{,}733{,}056$
exactly via Burnside's lemma (Davis's formula \cite{Davis54}),
independently of both nauty and OEIS;
re-verified all $9{,}924{,}592$ witness lines with its own sweep
verifier (zero failures); re-solved both minimizers with its own exact
solver, of a different design and not SAT-based, obtaining $9$ and $12$;
re-ran CaDiCaL afresh on the shipped CNFs; and exercised both stdlib
verifiers with negative controls --- deliberately corrupted artifacts
--- all of which were correctly rejected. The referee also imposed the
credit framing of \S\ref{ssec:credit} and caught one error in an
internal draft (Remark~\ref{rem:burnside}). On this protocol the result
was confirmed on 2026-08-05.

\section{The trusted base}\label{sec:tcb}

What a skeptic must believe, split into what is re-checked by short
readable code and what is assumed.

\emph{Machine-checked:} every witness packing, per isomorphism class,
from raw arc data (\S\ref{ssec:stdlib}); the totals and per-slice counts;
pairwise distinctness of the enumerated strings; the optimal packings in
both minimizers; the clause-multiset equality of each shipped CNF with an
independent regeneration from the specification; and the LRAT replay of
both unsatisfiability proofs, in a checker that refuses RAT steps.

\emph{Assumed:}
\begin{enumerate}
\item \emph{\texttt{gentourng} enumerates completely} --- one
representative of every isomorphism class. This is the one substantive
trust assumption, and it is supported rather than proved: exhaustive
cross-validation for $n\le7$; class counts at $n=9,10$ equal to OEIS
A000568 \cite{OEIS} and to the referee's independent exact Burnside
computation \cite{Davis54}; output strings pairwise distinct; and a
fresh re-run by the referee, equal to the shipped enumeration as a set
of strings byte for byte. Distinct strings with the right count do
not by themselves imply completeness (two listed strings could in
principle be isomorphic, hiding a missed class), which is why this item
is listed here and not above.
\item \emph{The sequential-counter encoding is complete} \cite{Sinz05}:
every packing of size $k$ extends to a satisfying assignment of the
counter variables. If the counter accidentally over-constrained, UNSAT
would mean less than claimed. (\cite{Sinz05} states the property; the
short conference paper omits the proofs for space.)
\item \emph{The arc-string convention is read as written.} Searcher and
verifiers necessarily share the \emph{convention} (not code) for
interpreting \texttt{gentourng}'s output format. A consistent
misreading on both sides would not be caught by line re-checking; it is
caught instead by the $n\le8$ reproduction of Yuster's values and by the
minimizers, whose blow-up structure is known independently of the
format.
\item CPython and the operating system.
\end{enumerate}

\noindent The searcher \texttt{tt3pack}, the encoder, CaDiCaL, and the
shipped CNF files are \emph{not} in the trusted base: every claim they
produced was re-derived from raw data or regenerated and re-checked
independently.

\section{The frontier: $n=11$}\label{sec:frontier}

The next open case is $\nut(11)$, conjectured value
$\lceil 110/6-11/3\rceil=15$. There are
$A000568(11)=903{,}753{,}248$ isomorphism classes, a factor of
about $93$ over the $n=10$ sweep --- CPU-days, not CPU-years, for the
pipeline exactly as shipped (\texttt{tt3pack} as written handles
$n\le11$; the $55$ arc pairs fit one $64$-bit mask). The corresponding
minimizer candidates are the blow-ups with parts $(4,4,3)$, whose
$6+6+3=15$ within-class arcs match the conjectured value, so both halves
are within reach of the same two-sided method. This
computation is the program's next target; it has not been started, and
no result at $n=11$ is claimed here. The fractional side ($\nuf(r)$ for $r=11$, which would
sharpen the $41/300$ of \cite{KY08}) is LP-based, separate, and was not
touched.

\begin{remark}[A corrected estimate]\label{rem:burnside}
An earlier internal draft of the working notes misstated $A000568(11)$
as ${\approx}9.04\times10^{11}$ and priced the $n=11$ sweep accordingly
(``${\sim}100$\,TB of logs''). The referee's exact Burnside computation
gives $903{,}753{,}248\approx9.04\times10^{8}$ --- the same mantissa,
three orders of magnitude smaller --- and the feasibility claim above
rests on the corrected value. We record the error because estimates that
die in private make the surviving ones less trustworthy.
\end{remark}

\section{Artifacts}\label{sec:artifacts}

The artifacts ship in two directories of the public repository
(\texttt{github.com/05oz/certify}): \texttt{tt3-certificates/} (sweeps,
certificates, minimizers --- $148$\,MB of gzipped $n=10$ sweep logs in
sixteen files, $2.3$\,MB for $n=9$) and \texttt{tt3-scripts/} (code and
verifiers), everything besides the sweep logs under $2$\,MB. The
certified release is exactly the $32$ files pinned by MD5
digest in \texttt{MD5SUMS.txt} and reproduced in Table~\ref{tab:md5}:
the sweeps, certificates, minimizers, verifiers, and code. Of the
auxiliary files cited in the text, \texttt{n10\_chunk\_counts.txt} and
\texttt{smalln\_table.txt} ship alongside them, unpinned;
\texttt{cand9.out}, \texttt{cand10.out}, and the archived primary-source
PDFs are retained in the author's private working tree and are not
redistributed (the PDFs for copyright reasons). The certified chain of
Theorem~\ref{thm:main} runs entirely through the pinned files.

To replay from scratch on any machine with nauty, CaDiCaL, a C compiler
and Python 3 (the commands are written for a flat working directory ---
in the repository the pinned files sit in \texttt{tt3-scripts/} and
\texttt{tt3-certificates/}; copy them into one directory first):

\begin{verbatim}
cc -O2 -o tt3pack tt3pack.c
gentourng -q 9  | ./tt3pack 9 9      # lower-bound sweep, n=9
gentourng -q 10 | ./tt3pack 10 12    # lower-bound sweep, n=10
python3 encode_cnf.py 9  "$(cat minimizer9.bits)"  10 min9_ge10.cnf
cadical --lrat --no-binary min9_ge10.cnf min9_ge10.lrat  # s UNSATISFIABLE
python3 verify_sweep.py --n 9 --target 9 --expect 191536 n9_sweep.txt.gz
python3 verify_minimizer.py --n 9 --value 9 --bits minimizer9.bits \
    --witness minimizer9.witness --cnf min9_ge10.cnf --lrat min9_ge10.lrat
\end{verbatim}

\noindent and the same pattern at $n=10$ with target $12$, value $12$,
and $K=13$ files. (The sweep commands stream to standard output;
redirect to a file. The first field of each sweep line is a per-run line
number, which restarts in every slice: when comparing an unsliced
$n=10$ replay against the sixteen shipped slices, compare the remaining
fields. Production ran the sixteen-slice form.) The two verifiers need only the Python standard
library; re-checking the shipped artifacts requires neither nauty, nor
CaDiCaL, nor a compiler, except for the enumeration itself, which is the
declared trust assumption of \S\ref{sec:tcb}. Tool versions used: nauty
2.9.3, CaDiCaL 3.0.1, Apple clang 17.0.0, CPython 3.

\begin{table}[ht]
\scriptsize
\setlength{\tabcolsep}{1pt}
\begin{center}
\begin{tabular}{ll@{\hspace{7pt}}ll}
\toprule
file & MD5 & file & MD5 \\
\midrule
\texttt{tt3pack.c} & \texttt{46592e2c895dcd88f485f35f408940dc} &
\texttt{n10\_sweep\_2.txt.gz} & \texttt{f9c71121e500d9594fab895a8fb9d6ba} \\
\texttt{encode\_cnf.py} & \texttt{c277f53e7343a44c93375159c248f08a} &
\texttt{n10\_sweep\_3.txt.gz} & \texttt{4e243e0e91582866523d512669c9eba4} \\
\texttt{verify\_sweep.py} & \texttt{a3726bfbb524e1d5d29bdfa892776fe0} &
\texttt{n10\_sweep\_4.txt.gz} & \texttt{d3d96b7b0a78432cb8e80dec1366a36a} \\
\texttt{verify\_minimizer.py} & \texttt{373639d979b8f44f559ca9ea78ef880b} &
\texttt{n10\_sweep\_5.txt.gz} & \texttt{bbe411ec67ba2e144ddd1fdb13cb1ce4} \\
\texttt{make\_candidates.py} & \texttt{546aabe9b6d4a28499da6be2c27e3b45} &
\texttt{n10\_sweep\_6.txt.gz} & \texttt{9fb5d3da724fcf453451e1ce4e639573} \\
\texttt{minimizer9.bits} & \texttt{34ed2059bf5e33ff439e040d0306b7a4} &
\texttt{n10\_sweep\_7.txt.gz} & \texttt{ece12815439e764a813c99a70b20ddb8} \\
\texttt{minimizer10.bits} & \texttt{66d67402bb479a6bc831d10d632c980c} &
\texttt{n10\_sweep\_8.txt.gz} & \texttt{b9fcb0b5939650e5b5a1f1fa51173438} \\
\texttt{minimizer9.witness} & \texttt{35a3d4136122454cb41624bd013c083c} &
\texttt{n10\_sweep\_9.txt.gz} & \texttt{1cd91ac37afb04e120415e5790a34bb6} \\
\texttt{minimizer10.witness} & \texttt{e1b9923118dbe2358007edf31da8da1f} &
\texttt{n10\_sweep\_10.txt.gz} & \texttt{a1e999d85ba88c3dd131fb123702b46b} \\
\texttt{min9\_ge10.cnf} & \texttt{f0080051d89ec538e86501efeeaecadf} &
\texttt{n10\_sweep\_11.txt.gz} & \texttt{64afda1fb8d36eebf04d70a53d2664b1} \\
\texttt{min9\_ge10.lrat} & \texttt{d13133b094e5bc0ab5fd7f7b7cf1af3b} &
\texttt{n10\_sweep\_12.txt.gz} & \texttt{90485a8cf9e63e716980c0ef74a6726c} \\
\texttt{min10\_ge13.cnf} & \texttt{ce16d4fc654198dcb1e3ff0782dd5b1a} &
\texttt{n10\_sweep\_13.txt.gz} & \texttt{0109f71fbb2a111405b175b6c40f2120} \\
\texttt{min10\_ge13.lrat} & \texttt{ad546eece47db3d3f7abf9baacc859bf} &
\texttt{n10\_sweep\_14.txt.gz} & \texttt{a1bb606d782c973e784672c6b90a3f11} \\
\texttt{min9\_ge9.cnf} & \texttt{a9acac4e212214975eb28f3690d21f86} &
\texttt{n10\_sweep\_15.txt.gz} & \texttt{a91a473a46034db7d448cd2eb6165019} \\
\texttt{min10\_ge12.cnf} & \texttt{eb5ae68eb5262236e2c0c8e4947b9c25} &
\texttt{n9\_sweep.txt.gz} & \texttt{b43c323b55c2f199d94cacc025d3d15b} \\
\texttt{n10\_sweep\_0.txt.gz} & \texttt{5069371e1500dda7eabdf5fe7a0d4e35} &
& \\
\texttt{n10\_sweep\_1.txt.gz} & \texttt{73e36c3db19e0815e57de78761d92dd2} &
& \\
\bottomrule
\end{tabular}
\end{center}
\medskip
\caption{MD5 digests of the released artifacts (the $32$ entries of
\texttt{MD5SUMS.txt}).}
\label{tab:md5}
\end{table}

\section*{Acknowledgments}

The problem, the conjectured formula, the verified range $n\le8$, and
the blow-up upper-bound argument used here at both orders are Raphael
Yuster's; the fractional method that still holds the best general
lower bound is Mohamad Kabiya and Raphael Yuster's. The enumeration
rests entirely on Brendan McKay and Adolfo Piperno's nauty. To Carsten
Sinz for the cardinality encoding, to Cruz-Filipe, Heule, Hunt, Kaufmann
and Schneider-Kamp for LRAT, and to Armin Biere and coauthors for
CaDiCaL. The computation and drafting were AI-assisted as stated in the
first footnote; the adversarial referee's insistence on the credit
framing of \S\ref{ssec:credit} made this a better and a more honest note.

\begin{thebibliography}{99}

\bibitem[KY08]{KY08}
M.~Kabiya and R.~Yuster,
\emph{Packing transitive triples in a tournament},
Ann.\ Comb.\ \textbf{12} (2008), no.~3, 291--306.
DOI 10.1007/s00026-008-0352-3. Read in full on August 5, 2026, from an
archived copy (\texttt{primary-source-kabiya-yuster-tt3.pdf}, retained
in the author's private working tree; not redistributed for copyright
reasons); that copy is the author
preprint, and section references follow its numbering: the blow-up
upper-bound argument is \S2.4, and the $r=10$ fractional LP computation
is \S3.

\bibitem[MP14]{nauty}
B.~D. McKay and A.~Piperno,
\emph{Practical graph isomorphism, II},
J.\ Symbolic Comput.\ \textbf{60} (2014), 94--112. Version used here:
nauty 2.9.3 (\texttt{gentourng}).

\bibitem[Dav54]{Davis54}
R.~L. Davis,
\emph{Structures of dominance relations},
Bull.\ Math.\ Biophys.\ \textbf{16} (1954), 131--140.

\bibitem[OEIS]{OEIS}
OEIS Foundation Inc.,
\emph{The On-Line Encyclopedia of Integer Sequences}, sequence A000568
(the number of tournaments on $n$ unlabeled nodes),
\url{https://oeis.org/A000568}.

\bibitem[OPG]{OPG}
Open Problem Garden, \emph{Minimum number of arc-disjoint transitive
subtournaments of order 3 in a tournament},
\url{http://www.openproblemgarden.org/op/minimum_number_of_transitive_subtournaments_of_order_3_in_a_tournament}.
Fetched live August 5, 2026; the entry records no progress beyond
\cite{KY08}.

\bibitem[BFF{+}24]{CaDiCaL}
A.~Biere, T.~Faller, K.~Fazekas, M.~Fleury, N.~Froleyks and F.~Pollitt,
\emph{CaDiCaL 2.0},
in: Computer Aided Verification --- CAV 2024, Lecture Notes in Comput.\
Sci.\ \textbf{14681}, Springer, 2024, 133--152.
DOI 10.1007/978-3-031-65627-9\_7. Version used to produce the proofs:
CaDiCaL 3.0.1 (\texttt{--lrat --no-binary}).

\bibitem[CFHKS17]{LRAT}
L.~Cruz-Filipe, M.~J.~H. Heule, W.~A. Hunt~Jr., M.~Kaufmann and
P.~Schneider-Kamp,
\emph{Efficient certified RAT verification},
in: Automated Deduction --- CADE-26, Lecture Notes in Comput.\ Sci.\
\textbf{10395}, Springer, 2017, 220--236.

\bibitem[Sin05]{Sinz05}
C.~Sinz,
\emph{Towards an optimal CNF encoding of Boolean cardinality
constraints},
in: Principles and Practice of Constraint Programming --- CP 2005,
Lecture Notes in Comput.\ Sci.\ \textbf{3709}, Springer, 2005, 827--831.

\bibitem[Yus04]{Yuster04}
R.~Yuster,
\emph{The number of edge-disjoint transitive triples in a tournament},
Discrete Math.\ \textbf{287} (2004), 187--191. Also
arXiv:math/0304180. The concluding remarks (\S4: the oriented Tur\'an
upper-bound construction; the verified range $n\le8$, with $n=8$ by
computer) were checked verbatim against the arXiv text on August 5,
2026 (archived copy retained in the author's private working tree; not
redistributed).

\end{thebibliography}

\end{document}
